% ============================================================================
% CH — EVALUATION FRAMEWORK                         [slug: eval-framework]
% STATUS: design frozen (S0 charter); reference numbers PENDING S0.7 fleet.
% Shared "evaluation design" chapter serving BOTH blocks (front matter,
% outside the \part{} markers).
% RESTRUCTURED 2026-07-14 (Matteo; granularity principle — headers follow
% content mass): 6 sections -> 2 + a reserved Part II placeholder. Old
% section labels parked on their new homes. Sub-headers only at prose time
% if blocks demand dividers.
% Material map: ../outline.md#eval-framework
% ============================================================================
\cleardoublepage
\chapter{Evaluation Framework (6--7)}   % page goal — scaffolding, DELETE before submission
\label{ch:eval}

\section{Computational Setup}
\label{sec:eval-compute}
\label{sec:eval-hyperopt}      % old hyperopt/infra section parked here
% Reproducibility corner (Matteo, 2026-07-14). Content blocks:
%  1. Hardware & usage: LEO5 HPC cluster (SLURM; A100/A40/A30) for all
%     training and hyperopt; laptop for development and small tests.
%  2. Hyperopt: Ray Tune search spaces (pointer to appendix for full
%     spaces); trial budgets — 40 configs for dev profiles, 100 for the
%     final runs (matching the ProtoMF paper's budget); early stopping
%     (patience 10); multi-seed protocol.
%  2b. TWO-TIER NUMBERS POLICY (Matteo, 2026-07-14): dev-profile runs
%     (incl. the S0.7 reference fleet) are the WORKING numbers during
%     development — comparison-valid because every model gets identical
%     treatment (C1-C8) — but explicitly NOT thesis numbers. Thesis-grade
%     tier (paper-budget hyperopt, 3 seeds averaged, mirroring ProtoMF's
%     methodology) is queued bit by bit ONLY after model freeze (no further
%     model changes). Cost adjustments only if the dev tier signals the
%     full tier is too expensive — decided then, not now. Every number
%     quoted in the thesis states its tier until the final swap.
%  3. Tracking: Weights & Biases logging; run manifests (hooks into charter
%     C8); configs versioned.
% Brief — details to appendix.

\section{Comparison Design}
\label{sec:eval-design}
\label{sec:eval-metrics}       % old section labels parked here —
\label{sec:eval-fleet}         % their content now lives in this
\label{sec:eval-charter}       % section as blocks
\label{sec:eval-replication}
\label{sec:eval-cold}
\vault{2026-06-16_1704_eval-design-citable-justification, 2026-06-16_1631_protomf-item-ranking-basecase, 2026-06-07_2308_map12-noncomparable-focus-shift, 2026-06-07_2309_eval-scope-lightweight-map12}
\vault{2026-07-11_1629_baseline-suite-published-rosters, 2026-06-16_1824_feature-aware-baselines, 2026-07-11_1556_bias-channel-popularity-prototype-interpretability, 2026-06-07_2306_replication-complete, 2026-03-09_1641_phase2-data-config}
\vault{2026-07-11_2023_coldstart-methodology-b5-anchoring, 2026-07-11_1526_facet-a-claim-boundary-coldstart-2x2, 2026-07-12_1659_five-core-loo-train-floor-and-label-hallucination-lesson}
% What is compared, how fairly, against whom. Content blocks:
%  1. Task + metrics: leave-one-out HR@10 / NDCG@10 (validation metric
%     hit_ratio@10); citable justification for position-aware metrics and
%     for NOT reporting RMSE/AUC; why MAP@12 leaderboard numbers are
%     structurally non-comparable (qualitative note only).
%  2. Comparison-fairness charter: the C1-C8 conditions (S0.6) — identical
%     splits, budgets, seeds, manifests; one run = one manifest. Likely a
%     table + short prose.
%  3. Reference fleet AND replication together (Matteo, 2026-07-14: "the
%     same thing with extension" — one paragraph + table): the fleet IS the
%     replicated ProtoMF roster (mf, acf, U-, I-, UI-ProtoMF), extended with
%     LightFM (tags / tags+ids) and the popularity floor; exclusion
%     reasoning for others; replication on the paper datasets (ml-1m,
%     amazon2014; lfm2b unavailable — license) establishes trust in the
%     codebase before extension. Replicated-vs-paper numbers table (main
%     text or appendix — decide at prose time).
%  4. Cold-item evaluation protocol (block, NOT shrunk to a mention —
%     the claim boundary does protective work): anchored on LightFM's
%     protocol (Kula 2015 §5), 20% item holdout, cold-pool ranking; five
%     declared deviations; cold-user 2x2 claim boundary (cold users
%     structurally out of reach for item-side designs — honest scope
%     statement). Split construction lives in the Data chapter
%     (\ref{sec:data-splits}); protocol + deviations live here.

\section{The Explanation System}
\label{sec:eval-explanations}
% Decided 2026-07-14 (Matteo): the explanation machinery lives HERE as shared
% infrastructure (parallel to the comparison design), demonstrated ON THE
% HOST ProtoMF (already introduced in Background) — the lineage chapters then
% show ONLY their extended capability ("focus on what is new"), never
% re-demonstrating the host's baseline explanations. No duplicates.
% Content blocks:
%  1. What an explanation IS in this family (Background supplies the
%     vocabulary): the additive score decomposition over prototypes —
%     per-prototype attribution of a recommendation, user- and item-side
%     read-outs.
%  2. The system/pipeline: attribution breakdown -> prototype naming ->
%     rendered artifacts (textual + graphical), generated automatically per
%     run. Naming has two modes: post-hoc neighborhood/centroid reading (all
%     the host can offer) vs parameter-derived naming through the
%     builder-owned vocabulary (what the grounded models unlock — described
%     generically here, demonstrated in the lineage chapters).
%  3. Demonstration on ProtoMF: one full recommendation explanation +
%     prototype interpretation walkthrough. Deliberately EXPOSES the naming
%     gap empirically — centroid-read prototype names, post-hoc and
%     unverifiable — the reader feels the problem right before the lineage
%     starts closing it. (Demo checkpoint: from the S0.7 fleet once landed.)
%  4. QA discipline: graphical outputs heavily scrutinized before landing
%     (standing directive); naming spot-checks.

\section{Evaluating Part II [PLACEHOLDER]}
\label{sec:eval-part2}
% RESERVED (Matteo, 2026-07-14): this chapter serves both blocks, so Part II
% ("Towards Interpretability as a Model Objective") gets its setup declared here
% too — "what we did there". Content deliberately unclear until the Part II
% experiments are designed (see iq-experiments): expected candidates —
% interpretability metrics/instruments (profile entropy/sharpness,
% feature-explained fraction, purity, determinacy), Rashomon accuracy
% tolerance (epsilon), knob-study protocol. Placeholder until then.
